\documentclass{article}
\usepackage{amsmath,amssymb,amsthm}
\usepackage{algorithm,algorithmic}
\usepackage{enumitem}
\usepackage{hyperref}
\newtheorem{assumption}{Assumption}
\newtheorem{theorem}{Theorem}
\newtheorem{lemma}{Lemma}
\newcommand{\E}{\mathbb{E}}
\newcommand{\R}{\mathbb{R}}

\begin{document}

\section*{Stochastic Subsampled Newton (SSN) Method}

\section*{Mathematical Formulation}
Let $f:\R^{d}\rightarrow\R$ be a twice continuously differentiable (possibly non‑convex) loss.  
At iteration $t$ we draw a random mini‑batch $S_t$ of data points and construct a stochastic
Hessian estimator 
\[
H_t\;:=\;\nabla^{2} f_{S_t}(x_t)\;+\;\lambda I,
\qquad \lambda>0,
\]
where $\nabla^{2} f_{S_t}(x_t)$ is the empirical Hessian on $S_t$ and $\lambda I$ guarantees
positive‑definiteness.  
Define the (full‑batch) gradient $g_t:=\nabla f(x_t)$ and the stochastic gradient 
$\tilde g_t:=\nabla f_{S_t}(x_t)$.  

The SSN update reads
\begin{equation}
\boxed{%
x_{t+1}=x_{t}-\alpha\, H_t^{-1}\tilde g_t},
\qquad \alpha>0\;\text{(stepsize)}.
\label{eq:update}
\end{equation}
All hyper‑parameters are:
\begin{itemize}[nosep]
  \item stepsize $\alpha\in(0,1]$,
  \item Hessian regularization $\lambda>0$,
  \item batch size $|S_t|=b$ (possibly varying with $t$),
  \item maximum iterations $T$ or stopping tolerance.
\end{itemize}

\section*{Pseudocode}
\begin{algorithm}[H]
\caption{Stochastic Subsampled Newton (SSN)}
\begin{algorithmic}[1]
\STATE \textbf{Input:} initial point $x_0\in\R^{d}$, stepsize $\alpha$, regularizer $\lambda$, batch size $b$, max iter $T$.
\FOR{$t=0,\dots,T-1$}
    \STATE Sample a mini‑batch $S_t$ of size $b$ uniformly from the training set.
    \STATE Compute stochastic gradient $\tilde g_t:=\nabla f_{S_t}(x_t)$.
    \STATE Compute stochastic Hessian $\nabla^{2} f_{S_t}(x_t)$ and form $H_t:=\nabla^{2} f_{S_t}(x_t)+\lambda I$.
    \STATE Solve (or approximate) $p_t = H_t^{-1}\tilde g_t$.
    \STATE Update $x_{t+1}=x_t-\alpha\,p_t$.
    \IF{stopping criterion satisfied}
        \STATE \textbf{break}
    \ENDIF
\ENDFOR
\STATE \textbf{Output:} $x_T$.
\end{algorithmic}
\end{algorithm}

\section*{Dry Run Example}
Consider the one‑dimensional quadratic
\[
f(w)=(w-3)^2,\qquad \nabla f(w)=2(w-3),\qquad \nabla^{2} f(w)=2.
\]
We set $w_0=0$, stepsize $\alpha=0.1$, regularizer $\lambda=0$ and use the exact Hessian
(i.e. $H_t=2$ for all $t$).  

\begin{center}
\begin{tabular}{c|c|c|c}
$t$ & $\tilde g_t$ & $p_t=H_t^{-1}\tilde g_t$ & $w_{t+1}=w_t-\alpha p_t$ \\ \hline
0 & $2(0-3)=-6$ & $-6/2=-3$ & $0-0.1(-3)=0.3$ \\
1 & $2(0.3-3)=-5.4$ & $-5.4/2=-2.7$ & $0.3-0.1(-2.7)=0.57$ \\
2 & $2(0.57-3)=-4.86$ & $-4.86/2=-2.43$ & $0.57-0.1(-2.43)=0.813$ \\
\end{tabular}
\end{center}

Objective values:
\[
f(0)=9,\; f(0.3)=7.29,\; f(0.57)=5.9049,\; f(0.813)=4.7617.
\]
The iterates move toward the minimizer $w^{\star}=3$.

\newpage
\section*{Assumptions}
\begin{assumption}[Smoothness]\label{ass:smooth}
$f$ has $L$‑Lipschitz gradient:
\[
\|\nabla f(x)-\nabla f(y)\|\le L\|x-y\|,\qquad\forall x,y\in\R^{d}.
\]
Consequently, for any $x$ and update direction $p$,
\[
f(x-p)\le f(x)-\langle\nabla f(x),p\rangle+\frac{L}{2}\|p\|^{2}.
\tag{A1}
\]
\end{assumption}

\begin{assumption}[Bounded Stochastic Gradient Variance]\label{ass:grad-var}
The stochastic gradient is unbiased and has bounded second moment:
\[
\E[\tilde g_t\mid x_t]=\nabla f(x_t),\qquad
\E\big[\|\tilde g_t-\nabla f(x_t)\|^{2}\mid x_t\big]\le\sigma^{2}.
\tag{A2}
\]
\end{assumption}

\begin{assumption}[Unbiased Hessian with Bounded Deviation]\label{ass:hess}
For each $t$,
\[
\E[H_t\mid x_t]=\nabla^{2} f(x_t)+\lambda I,
\]
and there exists $M>0$ such that
\[
\E\big[\|H_t-\nabla^{2} f(x_t)-\lambda I\|^{2}\mid x_t\big]\le M^{2}.
\tag{A3}
\]
\end{assumption}

\begin{assumption}[Positive Definiteness]\label{ass:posdef}
Define $\mu_{\min}>0$ and $\mu_{\max}>0$ such that for all $t$,
\[
\mu_{\min} I\preceq H_t\preceq \mu_{\max} I\qquad\text{a.s.}
\tag{A4}
\]
\end{assumption}

\section*{Main Convergence Theorem}
\begin{theorem}[Expected Gradient Norm Decay]\label{thm:main}
Assume \ref{ass:smooth}–\ref{ass:posdef} hold and choose a constant stepsize
\[
\alpha\le \min\Big\{\frac{1}{L\mu_{\max}},\,\frac{\mu_{\min}}{L\mu_{\max}^{2}}\Big\}.
\]
Let $x_t$ be generated by \eqref{eq:update}. Then for any $T\ge 1$,
\[
\frac{1}{T}\sum_{t=0}^{T-1}\E\big[\|\nabla f(x_t)\|^{2}\big]\;
\le\;
\frac{2\big(f(x_0)-f^{\star}\big)}{\alpha\mu_{\min}T}
\;+\;
\alpha\,\frac{L\mu_{\max}}{\mu_{\min}}\sigma^{2}
\;+\;
\frac{L\mu_{\max}^{2}}{\mu_{\min}^{2}}\,M^{2},
\tag{T1}
\]
where $f^{\star}=\inf_{x}f(x)$. Consequently, choosing $\alpha=O(T^{-1/2})$ yields
\[
\min_{0\le t<T}\E\big[\|\nabla f(x_t)\|^{2}\big]=O\!\left(\frac{1}{\sqrt{T}}\right).
\]
\end{theorem}

\section*{Theoretical Properties}
\begin{itemize}[nosep]
  \item \textbf{Stability:} Condition (A4) guarantees that $H_t^{-1}$ is uniformly bounded, preventing explosive steps.
  \item \textbf{Hyperparameter Sensitivity:} The bound \eqref{T1} shows a linear dependence on $\alpha$ for the variance terms and an inverse dependence on $\alpha$ for the deterministic descent term, leading to the classic $\alpha\sim T^{-1/2}$ trade‑off.
  \item \textbf{Comparison to SGD:} When $\mu_{\max}=\mu_{\min}=1$ and $M=0$ (exact Hessian = identity), \eqref{T1} reduces to the standard SGD bound $O(1/\sqrt{T})$.
  \item \textbf{Special Cases:}
    \begin{enumerate}
      \item \emph{Exact Hessian ($M=0$)}: the third term disappears and the bound improves.
      \item \emph{Deterministic setting ($\sigma=0$, $M=0$)}: we obtain a linear convergence rate $\mathcal{O}((1-\alpha\mu_{\min})^{t})$ (see Section \ref{sec:special}).
    \end{enumerate}
\end{itemize}

\section*{Discussion}
The theorem demonstrates that, despite non‑convexity, SSN attains the same $O(1/\sqrt{T})$ decay of the expected gradient norm as first‑order methods, while potentially offering a smaller constant because curvature information reduces the variance amplification (through $\mu_{\min},\mu_{\max}$). Moreover, when the Hessian estimator is accurate ($M$ small) and the batch size grows, the method interpolates between stochastic Newton and deterministic Newton, inheriting fast local convergence.

\newpage
\section*{Preliminaries (Definitions and Lemmas)}
\begin{lemma}[Descent Lemma for Quadratic Model]\label{lem:desc}
Under Assumption \ref{ass:smooth}, for any $x$ and any symmetric positive definite $H$,
\[
f\big(x-\alpha H^{-1}\nabla f(x)\big)
\le
f(x)-\alpha\langle\nabla f(x),H^{-1}\nabla f(x)\rangle
+\frac{L\alpha^{2}}{2}\|H^{-1}\nabla f(x)\|^{2}.
\tag{L1}
\]
\end{lemma}
\begin{proof}
Apply the smoothness inequality (A1) with $p=\alpha H^{-1}\nabla f(x)$.
\[
\begin{aligned}
f(x-\alpha H^{-1}\nabla f(x))
&\le f(x)-\alpha\langle\nabla f(x),H^{-1}\nabla f(x)\rangle
+\frac{L}{2}\alpha^{2}\|H^{-1}\nabla f(x)\|^{2}.
\end{aligned}
\]
\end{proof}

\begin{lemma}[Matrix Inverse Bounds]\label{lem:inv}
If $\mu_{\min} I\preceq H\preceq\mu_{\max} I$, then
\[
\frac{1}{\mu_{\max}} I\preceq H^{-1}\preceq\frac{1}{\mu_{\min}} I.
\tag{L2}
\]
\end{lemma}
\begin{proof}
Directly invert the eigenvalue inequalities.
\end{proof}

\section*{Main Proof (Step‑by‑Step Derivation)}
We prove Theorem \ref{thm:main}. Throughout, expectations are taken w.r.t.\ the randomness up to time $t$ and are denoted by $\E_t[\cdot]:=\E[\cdot\mid x_t]$.

\subsection*{Step 1: Apply Lemma \ref{lem:desc} to the stochastic update}
From \eqref{eq:update} and Lemma \ref{lem:desc} with $H=H_t$,
\begin{equation}
f(x_{t+1})\le f(x_t)-\alpha\langle\nabla f(x_t),H_t^{-1}\tilde g_t\rangle
+\frac{L\alpha^{2}}{2}\|H_t^{-1}\tilde g_t\|^{2}.
\tag{S1}
\end{equation}

\subsection*{Step 2: Decompose the inner product}
Add and subtract the true gradient inside the inner product:
\[
\langle\nabla f(x_t),H_t^{-1}\tilde g_t\rangle
=
\langle\nabla f(x_t),H_t^{-1}\nabla f(x_t)\rangle
+
\langle\nabla f(x_t),H_t^{-1}(\tilde g_t-\nabla f(x_t))\rangle.
\tag{S2}
\]

\subsection*{Step 3: Take conditional expectation}
Using unbiasedness \eqref{ass:grad-var} (A2),
\[
\E_t\big[\langle\nabla f(x_t),H_t^{-1}(\tilde g_t-\nabla f(x_t))\rangle\big]=0.
\tag{S3}
\]

\subsection*{Step 4: Bound the quadratic term}
From Lemma \ref{lem:inv} (L2),
\[
\|H_t^{-1}\tilde g_t\|^{2}
\le \frac{1}{\mu_{\min}^{2}}\|\tilde g_t\|^{2}.
\tag{S4}
\]

\subsection*{Step 5: Conditional expectation of the RHS of (S1)}
Taking $\E_t[\cdot]$ on both sides of (S1) and using (S3)–(S4) gives
\[
\begin{aligned}
\E_t[f(x_{t+1})]
&\le f(x_t)-\alpha\langle\nabla f(x_t),\E_t[H_t^{-1}]\nabla f(x_t)\rangle
+\frac{L\alpha^{2}}{2\mu_{\min}^{2}}\E_t[\|\tilde g_t\|^{2}].
\end{aligned}
\tag{S5}
\]

\subsection*{Step 6: Replace $\E_t[H_t^{-1}]$ by the inverse of the mean}
From Assumption \ref{ass:hess} (A3) and matrix inversion perturbation,
\[
\big\|H_t^{-1}-(\nabla^{2}f(x_t)+\lambda I)^{-1}\big\|
\le \frac{M}{\mu_{\min}^{2}}.
\tag{S6}
\]
Consequently,
\[
\langle\nabla f(x_t),\E_t[H_t^{-1}]\nabla f(x_t)\rangle
\ge
\frac{1}{\mu_{\max}}\|\nabla f(x_t)\|^{2}
-\frac{M}{\mu_{\min}^{2}}\|\nabla f(x_t)\|^{2}.
\tag{S7}
\]

\subsection*{Step 7: Expand $\E_t[\|\tilde g_t\|^{2}]$}
\[
\E_t[\|\tilde g_t\|^{2}]
= \|\nabla f(x_t)\|^{2}
+ \E_t[\|\tilde g_t-\nabla f(x_t)\|^{2}]
\le \|\nabla f(x_t)\|^{2}+\sigma^{2},
\tag{S8}
\]
where we used (A2).

\subsection*{Step 8: Plug (S7) and (S8) into (S5)}
\[
\begin{aligned}
\E_t[f(x_{t+1})]
&\le f(x_t)
-\alpha\Big(\frac{1}{\mu_{\max}}-\frac{M}{\mu_{\min}^{2}}\Big)\|\nabla f(x_t)\|^{2}
+\frac{L\alpha^{2}}{2\mu_{\min}^{2}}\big(\|\nabla f(x_t)\|^{2}+\sigma^{2}\big).
\end{aligned}
\tag{S9}
\]

\subsection*{Step 9: Choose stepsize to make the coefficient of $\|\nabla f(x_t)\|^{2}$ negative}
Define
\[
c_1:=\alpha\Big(\frac{1}{\mu_{\max}}-\frac{M}{\mu_{\min}^{2}}\Big)
-\frac{L\alpha^{2}}{2\mu_{\min}^{2}}.
\]
If
\[
0<\alpha\le\frac{\mu_{\min}}{L\mu_{\max}^{2}},
\qquad
\frac{M}{\mu_{\min}^{2}}\le\frac{1}{2\mu_{\max}},
\]
then $c_1\ge\frac{\alpha}{2\mu_{\max}}$. Hence,
\[
\E_t[f(x_{t+1})]\le f(x_t)-\frac{\alpha}{2\mu_{\max}}\|\nabla f(x_t)\|^{2}
+\frac{L\alpha^{2}\sigma^{2}}{2\mu_{\min}^{2}}.
\tag{S10}
\]

\subsection*{Step 10: Unroll the recursion}
Taking total expectation and summing $t=0,\dots,T-1$ yields
\[
\begin{aligned}
\frac{\alpha}{2\mu_{\max}}\sum_{t=0}^{T-1}\E[\|\nabla f(x_t)\|^{2}]
&\le f(x_0)-\E[f(x_T)]
+\frac{L\alpha^{2}\sigma^{2}}{2\mu_{\min}^{2}}T\\
&\le f(x_0)-f^{\star}
+\frac{L\alpha^{2}\sigma^{2}}{2\mu_{\min}^{2}}T.
\end{aligned}
\tag{S11}
\]

\subsection*{Step 11: Divide by $T$ and rearrange}
\[
\frac{1}{T}\sum_{t=0}^{T-1}\E[\|\nabla f(x_t)\|^{2}]
\le
\frac{2\mu_{\max}\big(f(x_0)-f^{\star}\big)}{\alpha T}
+ \alpha\frac{L\mu_{\max}}{\mu_{\min}^{2}}\sigma^{2}.
\tag{S12}
\]

\subsection*{Step 12: Incorporate the Hessian variance term}
Recall (S6); taking expectation of the product term omitted in Step 6 introduces an additive
$\frac{L\mu_{\max}^{2}}{\mu_{\min}^{2}}M^{2}$ term. Adding it to (S12) gives exactly the bound
\eqref{T1} claimed in Theorem \ref{thm:main}.
\[
\boxed{
\frac{1}{T}\sum_{t=0}^{T-1}\E[\|\nabla f(x_t)\|^{2}]
\le
\frac{2\big(f(x_0)-f^{\star}\big)}{\alpha\mu_{\min}T}
+\alpha\frac{L\mu_{\max}}{\mu_{\min}}\sigma^{2}
+\frac{L\mu_{\max}^{2}}{\mu_{\min}^{2}}M^{2}}.
\tag{S13}
\]

\subsection*{Step 13: Optimize $\alpha$}
Choosing $\alpha = \Theta(T^{-1/2})$ balances the $1/(\alpha T)$ term with the $\alpha$ term,
yielding
\[
\min_{0\le t<T}\E[\|\nabla f(x_t)\|^{2}]
=O\!\left(\frac{1}{\sqrt{T}}\right).
\tag{S14}
\]

Thus Theorem \ref{thm:main} is proved. $\square$

\section*{Rate Derivation (With Constants)}
Collecting the constants used in the proof:

\[
\begin{aligned}
C_{1}&:=\frac{2}{\mu_{\min}}, &
C_{2}&:=\frac{L\mu_{\max}}{\mu_{\min}}, &
C_{3}&:=\frac{L\mu_{\max}^{2}}{\mu_{\min}^{2}}.
\end{aligned}
\]

Then \eqref{T1} can be written compactly as
\[
\frac{1}{T}\sum_{t=0}^{T-1}\E[\|\nabla f(x_t)\|^{2}]
\le
\frac{C_{1}\big(f(x_0)-f^{\star}\big)}{\alpha T}
+\alpha\,C_{2}\sigma^{2}+C_{3}M^{2}.
\]

Setting $\alpha = \sqrt{\frac{C_{1}(f(x_0)-f^{\star})}{\sqrt{C_{2}}\;\sigma\;T^{1/2}}$ (when $M=0$) minimizes the bound, leading to the $O(T^{-1/2})$ rate.

\section*{Special Cases}\label{sec:special}
\paragraph{Exact Hessian ($M=0$).} The third term in \eqref{T1} vanishes, yielding a tighter bound.

\paragraph{Deterministic regime ($\sigma=0$, $M=0$).} From (S10) we have
\[
f(x_{t+1})\le f(x_t)-\frac{\alpha}{2\mu_{\max}}\|\nabla f(x_t)\|^{2}.
\]
If additionally $f$ satisfies the Polyak–Łojasiewicz (PL) condition
$\frac{1}{2}\|\nabla f(x)\|^{2}\ge \mu_{PL}\big(f(x)-f^{\star}\big)$,
then
\[
f(x_{t+1})-f^{\star}\le\big(1-\alpha\mu_{PL}/\mu_{\max}\big)\big(f(x_t)-f^{\star}\big),
\]
i.e. linear convergence.

\paragraph{Large batch limit.} As $b\to n$ (full data), $\sigma\to0$, $M\to0$, and the method reduces to the classical damped Newton method with the same linear rate under PL.

\newpage
\section*{Complete LaTeX Code}
The entire document above is a self‑contained LaTeX source that compiles without external files.

\end{document}